\subsection{Análisis del Módulo 4}

Basado en la \cref{tab:framework}, se selecciona React para el Módulo 4, 
dado que su enfoque en la construcción de interfaces de usuario interactivas mediante componentes y 
JSX permite una integración eficiente con los datos JSON enviados por el módulo 1 a través de MQTT. Su alto 
rendimiento, gracias al DOM Virtual, asegura actualizaciones rápidas de la UI para visualizar datos dinámicos de 
fermentación en tiempo real, mientras que su amplio ecosistema de librerías facilita la incorporación de herramientas 
para enrutamiento y gestión de estado, adaptándose a los requerimientos funcionales \ref{it:RF07} y no funcionales RNF04. 
Aunque requiere librerías adicionales, lo que incrementa el tamaño, su comunidad activa y extensa documentación compensan 
esta limitación, garantizando compatibilidad con el backend Python y MongoDB mediante APIs ligeras. En contraste, 
Angular es más rígido por su estructura MVC y dependencia de TypeScript, mientras que no se descarta el uso Vue.js, 
ya que es muy ligero, 
tiene un ecosistema más reducido que limita las soluciones prediseñadas. Por consiguiente, React es la opción 
óptima para desarrollar una interfaz escalable y reactiva en este sistema IoT.

\begin{table}[htbp] % El entorno table es necesario para \caption y \label
    \centering
    \caption{Comparación de Frameworks}
    \label{tab:framework}
    \begingroup
    \setlength{\tabcolsep}{3pt} 
    \renewcommand{\arraystretch}{1.5} 
    \small 
    \begin{tabularx}{\textwidth}{|l|l|X|X|X|}
    \hline
    \textbf{Tipo} & \textbf{Nombre} & \textbf{Enfoque Principal} & \textbf{Ventaja} & \textbf{Desventaja} \\ \hline
    
    \textbf{Librería} & \textbf{React} \cite{zotero-item-134} & Construcción de Interfaces de Usuario (UI) basadas en componentes y JSX. & Alto rendimiento gracias al \textbf{DOM Virtual} y amplio ecosistema. & No es un framework completo; requiere librerías adicionales para enrutamiento. \\ \hline
    
    \textbf{Framework} & \textbf{Angular} \cite{zotero-item-136} & Marco MVC para aplicaciones a gran escala (SPA). & Solución \textit{full-featured} con TypeScript y arquitectura modular nativa. & Curva de aprendizaje alta, robusto y obligatorio el uso de TypeScript. \\ \hline
    
    \textbf{Framework} & \textbf{Vue.js} \cite{zotero-item-138} & Modelo declarativo basado en componentes con HTML, CSS y JS estándar. & Diseño progresivo, muy ligero y flexible para integraciones rápidas. & Ecosistema más pequeño; puede requerir desarrollo manual de funciones avanzadas. \\ \hline
    \end{tabularx}
    \endgroup
\end{table}

Basado en la \cref{tab:graficas}, se selecciona Recharts para la visualización de datos en el Módulo 4,
debido a su enfoque en componentes React, lo que facilita la integración directa.

\begin{table}[H] % [H] fuerza la posición exacta debajo de este texto
    \caption{Librerías de gráficos para React}
    \label{tab:graficas}
    \begingroup
    \setlength{\tabcolsep}{4pt} 
    \renewcommand{\arraystretch}{1.5} 
    \footnotesize 
    % Quitamos el \centering para que no se fuerce al centro si el ancho fuera menor
    \begin{tabularx}{\textwidth}{|>{\RaggedRight}p{1.9cm}|>{\RaggedRight}p{1.7cm}|>{\RaggedRight}X|>{\RaggedRight}p{2.2cm}|>{\RaggedRight}X|}
    \hline
    \textbf{Librería} & \textbf{Tipo de Gráficos} & \textbf{Flexibilidad} & \textbf{Rendimiento} & \textbf{Comentarios / Uso Principal} \\
    \hline
    \textbf{Recharts} \cite{zotero-item-220} & Compon- entes React / SVG & Moderada; ideal para gráficos estándar. & Alto & Solución declarativa, adecuada para métricas de fermentación en tiempo real. \\
    \hline
    \textbf{React-Chartjs-2} \cite{zotero-item-222} & Compon- entes React / Canvas & Media; basada en la API de Chart.js. & Medio-alto & Eficiente para series temporales de sensores y alta frecuencia. \\
    \hline
    \textbf{Nivo} \cite{zotero-item-224} & SVG, Canvas, HTML & Alta; modularidad y herramientas interactivas. & Alto & Ideal para dashboards dinámicos con datos de IoT profesionales. \\
    \hline
    \textbf{Victory} \cite{zotero-item-226} & Componentes React / SVG & Alta; API consistente para gráficos reutilizables. & Medio & Enfocada en dashboards para monitoreo de múltiples parámetros. \\
    \hline
    \end{tabularx}
    \endgroup
\end{table}

\subsection{Framework}

Basado en la tabla comparativa, se selecciona Flask para el Módulo 4, 
ya que su naturaleza como microframework proporciona flexibilidad y extensibilidad,
permitiendo una integración directa con Python y JSON para gestionar datos de fermentación
enviados por ESP32 vía MQTT, almacenados en MongoDB. Su estructura MVC adaptable optimiza 
el desarrollo de una aplicación pequeña en un solo servidor, facilitando la conexión con React 
mediante APIs ligeras. Aunque requiere configuración manual, esto es aceptable para un sistema compacto. 
En contraste, Django es rígido y orientado a SQL, mientras que FastAPI, enfocado en APIs complejas, excede 
las necesidades actuales. Por ello, Flask es la mejor opción por su simplicidad y escalabilidad.

\begingroup
\setlength{\tabcolsep}{3pt} 
\renewcommand{\arraystretch}{1.3} % Bajamos un poco el espacio para que quepa mejor
\footnotesize % Bajamos a footnote para asegurar que la tabla no sea tan alta

% Usamos longtable para que si NO cabe, se parta y no deje la hoja anterior vacía
\begin{longtable}{|p{1.9cm}|>{\RaggedRight}p{2.8cm}|>{\RaggedRight}p{2.8cm}|>{\RaggedRight}p{3.2cm}|>{\RaggedRight}p{3.2cm}|}
    \caption{Comparación de Frameworks para el módulo 4.} \label{tab:frameworks} \\
    \hline
    \textbf{Framework} & \textbf{Modularidad} & \textbf{Estructura} & \textbf{Ventaja} & \textbf{Desventaja} \\ \hline
    \endfirsthead
    
    \multicolumn{5}{l}{\footnotesize\itshape -- continuación de la página anterior} \\ \hline
    \textbf{Framework} & \textbf{Modularidad} & \textbf{Estructura} & \textbf{Ventaja} & \textbf{Desventaja} \\ \hline
    \endhead

    \textbf{Django} \cite{zotero-item-140} & Alta modularidad vía \textit{apps}; extensibilidad enfocada a SQL. & Patrón \textbf{MTV} rígido; framework "con opinión". & Solución \textit{full-stack} completa: ORM, admin y seguridad integrados. & Rigidez estructural; integración compleja con MongoDB. \\ \hline
    
    \textbf{Flask} \cite{zotero-item-142} & Microframework con excelente modularidad y máxima extensibilidad. & Patrón \textbf{MVC} flexible; el desarrollador define la organización. & Integración directa y de bajo acoplamiento con NoSQL (JSON). & Requiere configuración manual de cada librería externa. \\ \hline
    
    \textbf{FastAPI} \cite{zotero-item-144} & Alta modularidad mediante \textit{routers} y tipado estándar. & Basado en funciones de ruta y modelos \textit{Pydantic}. & Generación automática de documentación (\textit{Swagger}) y tipado rápido. & Enfocado estrictamente en APIs; no incluye motor de \textit{templates}. \\ \hline
\end{longtable}
\endgroup
